% 本书编写规范

\pentry{本书编写规范\upref{Sample}}
\subsection{软件使用规范}
本书使用 TeXLive2016 软件中的 XeLaTeX 进行编译（暂不兼容更高版本）．如果 Windows 中编译卡在 eu1lmr.fd 上的时间较长，说明 font config 有问题，在 Windows 的控制行运行“fc-cache -fv”，重启 TeXLive，多试几次即可．TeXWork 编辑器中 Ctrl+T 编译，Ctrl+单击跳转到对应的 pdf 或代码，在 pdf 中 Alt+左箭头返回上一个位置． 代码中\textbackslash beq+Tab 生成公式环境，\textbackslash sub+Tab 生成 subsection．Ctrl+F 进行查找，Ctrl+G 查找下一个．菜单中的 Edit>Preference 设置默认字体为 Microsoft YaHei UI（11pt）， 默认编译器为 XeLaTeX，编码选择 UTF-8．

搜索文件夹内所有文档的内容用 FileSeek 软件，搜索空格用“\textbackslash 空格”，搜索“\$”用“\textbackslash \$”，以此类推．对比两个文档或文件夹用 WinMerge 软件．

画图用 Adobe Illustrator 和 Autodesk Graphics，用 MathType 在图中添加公式，希腊字母粗体正体矢量用从 Symbol 字体中插入（勾选 bold），更简单的方法是，先输入希腊字母，选中，然后在 Style 里面选 Vector-Matrix． 图片中的文字必须是 12 号， 如果字太小， 就把图片缩小而不是字放大． 图片中的线条尽量用 1pt 粗细． 图片插入书中后， 图中的字体应比书中的略小． 要画箭头， 先画一条直线， 然后选画笔图标， 在左下角的菜单中选 Arrows>Arrow\_Standard， 选 1.23 号箭头， 再把直线改成 0.33pt （相当于 1pt 粗的直线）．

\subsection{文件版本管理}
使用 GitHub Desktop， 用 PhysWiki repository管理所有文件， 每次 commit 需要做的事情如下
\begin{itemize}
\item 用 FileSeek 替换所有文档中的空心句号．
\item 确保所有文档可以顺利编译．
\item 解决编译产生的 warning．
\item 检查变化的内容．
\item 把 ManicTime 记录的写作时间记录到“计时.txt”．
\end{itemize}
每次 commit 的标题必须是下列之一
\begin{itemize}
\item 常规更新：包括完善词条，新词条等．
\item 词条统计：统计文件夹，对照表，和书中的词条，查看不一致或缺失．
\item 模板更新：模板有更新．
\item 批量修改：在多个文件中修改某一格式规范， 这种修改比较危险， 需要谨慎．
\end{itemize}
Repository 中创建 “.gitignore” 文件，忽略所有以下后缀名的文件：*.toc， *.aux， *.log， *.out， *synctex.gz．

词条统计的方法： 首先把 contents 文件夹中的所有文件名按顺序排列， 复制到表格中， 然后把词条对照表中的所有标签在表格中找到对应项， 做标记， 并把对照表中的词条名粘贴到表格中． 最后到 PhysWiki.tex 中逐个把标签在表格中找到对应项， 做标记， 对照词条名， 并对照词条文件中第一行的词条名．

\subsection{词条编写规范}

每个词条文件必须有一个独一无二的标签， 词条标签必须限制在 6 个字符内，必须在 PhysWiki.tex， 词条标签对照表和词条文件名中一致． 词条的中文名必须在主文件， 词条标签对照表和词条文件的第一行注释中一致． 如果不是超纲词条， 就把词条放在 contents 目录下， 并在主文件中用 \textbackslash entry \{\}\{\} 命令输入中文名和标签． 超纲词条放在 contents1 目录下并用 \textbackslash Entry \{\}\{\} 命令． 主文件有 Debug.tex， PhysWiki1.tex， 和 PhysWiki.tex 三个， 他们共用一个模板（others 目录）． 新词条必须现在 Debug.tex 中编辑， 完成后再从中删除并将 entry/Entry 插入到 PhysWiki.tex（超纲词条） 或 PhysWiki1.tex 中（力学分册）． Debug.tex 中的 entry/Entry 命令的后面可以用 \textbackslash newpage 命令强制换页， 但不允许在 PhysWiki.tex/PhysWiki1.tex 中这么做．

引用词条用 \textbackslash upref 命令，“预备知识”用 \textbackslash pentry 命令，“应用实例”用 \textbackslash eentry 命令，拓展阅读用 \textbackslash rentry 命令．

\subsubsection{黑色的小标题}

正文必须使用中文的括号，逗号，引号，冒号，分号，问号，感叹号，以及英文的全角实心的句号． 如 “牛顿—莱布尼兹公式” 中的横线必须用中文单破折号． 所有的标点符号前面不能有空格，后面要有空格． 正文\bb{粗体}用 \textbackslash bb， 等效于 \textbackslash textbf． 行内公式用单个美元符号，且两边要有空格，例如 $a^2+b^2=c^2$， 后面有标点符号的除外．方便的办法是先全部使用中文标点，最后再把所有空心句号替换成全角实心句号．

公式的 label 必须要按照“词条标签\_eq编号”的格式，只有需要引用的公式才加标签，编号尽量与显示的编号一致，但原则上不重复即可．图表的标签分别把 eq 改成 fig 和 tab 即可，例题用 ex， 习题用 exe．但凡是有 \textbackslash caption 命令的，\textbackslash label 需要紧接其后． 公式严禁使用 MathType 的自动编码， 必须手打．
\begin{equation}\label{Sample_eq1}
(a+b)^n = \sum_{i=0}^n C_n^i a^i b^{n-i} \quad (n\text{为整数})
\end{equation}
引用公式和图表都统一使用 \textbackslash autoref 命令，注意前面不加空格后面要加空格（后面是标点符号除外）， 例如\autoref{Sample_eq1}． 如果要引用其他词条中的公式，可以引用“其他词条\upref{Sample}”的\autoref{Sample_eq1} 也可以用“\autoref{Sample_eq1}\upref{Sample}”， 为了方便在纸质书上使用， 词条页码是不能忽略的．

\subsection{错别字替换}
正文中常见的错别字如“他们”（它们）， “一下”（以下）， 可以时常搜索替换．

\subsection{公式规范}
公式中的空格从小到大如 $a\, b\; c\quad d\qquad e$， 微分符号如 $\dd{x}$， 注意大括号不可省略． 三点省略号用 $\dots$， 如果要强制在下方， 用 $\ldots$． 大于等于或小于等于必须用 $\les, \ges$ 代替 $\le, \ge$． 自然对数底如 $\E$， 双重极限如
\begin{equation}
\lim_{\substack{\Delta x_i\to 0\\ \Delta y_i\to 0}} \sum_{i, j} f(x_i,y_i) \Delta x_i \Delta y_j
\end{equation}
注意行内公式的 \textbackslash lim 和 \textbackslash sum 如果要在正上方或正下方写东西， 需要加 \textbackslash limits， 如 $\lim\limits_{x\to 0}$．

可变尺寸的小中大括号一律用 Physics 宏包里面的 \textbackslash qty， 如
\begin{equation}
\qty(\frac ab)^2 \qquad \qty[\frac ab]^2 \qquad \qty{ \frac ab}
\end{equation}
对易算符如 $\comm{\Q A}{\Q B}$ 或 $\comm*{\Q A}{\Q B}$， 前者自动尺寸， 后者强制小尺寸． 泊松括号如 $\pb{\frac12}{B}$ 和 $\pb*{\frac12}{B}$．

如果要强制分式立体显示， 用 \textbackslash dfrac 命令， 如果 \textbackslash frac 后面只有两个字符的代码，最好写成如 $\frac12, \frac ab, \frac1a$ 等． 斜分数线后面若多于一个变量需要加括号， 如 $ab/(cd)$． 另外， 常用三角函数和自然指数对数函数后面的小括号会自动改变尺寸（Physics 宏包功能），若要给他们加幂，用中括号
\begin{equation}
\sin(\frac ab) \qquad \sin[2](\frac ab)
\end{equation}
另外可以用自定义的 \textbackslash sinc 函数如 $\sinc x, \sinc(x), \sinc^2(x+y)$，但后面的括号没有自动尺寸．

导数和偏导用 Physics 宏包里面的
\begin{equation}
\dv{x} \quad \dv{f}{x} \quad \dv[2]{f}{x} \quad \dv*[2]{f}{x} \quad
\pdv{x} \quad \pdv{f}{x} \quad \pdv[2]{f}{x} \quad \pdv{f}{x}{y} \quad \pdv*[2]{f}{x}
\end{equation}
复数如 $u+\I v$，复共轭如 $z\Cj$ 实部和虚部如 $\Re[z], \Im[z]$，  行内分式如 $a/b$， 不允许行内用立体分式．公式中的绝对值如 $\abs{a}$， 定积分求值如 $\eval{x^2}_0^1$ （自动尺寸）， 矢量如 $\vec a$， 手写矢量如 $\overrightarrow{a}$， 单位矢量如 $\uvec a$， 矢量点乘如 $\vec A \vdot \vec B$ （不可省略），矢量叉乘如 $\vec A\cross\vec B$． 量子力学算符如 $\Q a$（一般可以不加， 只有必要的时候加）， 矢量量子算符如 $\Qv p$． 狄拉克符号（加 * 号强制小尺寸）如 $\bra{a}, \ket{b}, \braket{a}, \braket{a}{b}, \mel{a}{\Q Q}{b}$． 平均值（加 * 号强制小尺寸）如 $\ev*{\Q Q}, \ev{\Q Q}{\psi}$． 梯度散度旋度拉普拉斯如 $\vec\nabla V$，$\vec\nabla\vdot\vec A$，$\vec\nabla\cross\vec A$，  $\nabla^2 V$， 但最好用 $\grad V$，$\div\vec A$，$\curl\vec A$，  $\laplacian V$． 单独一个粗体的 $\nabla$ 用 $\Nabla$． 行列式，矩阵 $\mat A$，转置 $\mat A \Tr$，厄米共轭 $\mat A \Her$ 如
\begin{equation}
\vmat{1&2\\3&4} \qquad 
\pmat{1&2\\3&4} \qquad 
\pmat{1&2\\3&4}\Tr \qquad
\pmat{1&2\\3&4}\Her
\end{equation}
行内的列矢量用行矢量的转置表示，如 $(1,2,3)\Tr$． 张量如 $\ten T$．

行间公式换行及对齐用 aligned 环境， 或用自定义的 \textbackslash ali 命令， 注意该环境可嵌套．
\begin{equation}\begin{aligned}
(a-b)^2 &= a^2+b^2 - 2ab \\
& = a^2+b^2+2ab-4ab\\
& = (a+b)^2-4ab
\end{aligned}\end{equation}
\begin{equation}\ali{
k_1 &= f(y_n, t_n) 
& k_2 &= f \qty(y_n + h\frac{k_1}{2}, t_n + \frac h2 )\\
k_3 &= f \qty( y_n + h\frac{k_2}{2}, t_n + \frac h2 ) \qquad
& k_4 &= f(y_n + hk_3, t_n + h)
}\end{equation}
左大括号用自定义的 \textbackslash leftgroup 命令， 里面相当于 aligned 环境
\begin{equation}
\leftgroup{
d+e+f&=g\\
a+b&=c}
\end{equation}
可变化尺寸的斜分数线如下
\begin{equation}
\left. \dv[2]{X}{x} \middle/ X + \dv[2]{Y}{y} \middle/ Y + \dv[2]{Z}{z} \middle/ Z  = \frac{1}{c^2}  \dv[2]{T}{t} \middle/ T\right.
\end{equation}
希腊字母如下
\begin{equation}\ali{
&\alpha (a), \beta (b), \chi (c), \delta (d), \epsilon/\varepsilon (e), \phi (f), \gamma (g),
\eta (h), \iota (i), \varphi (j), \kappa (k), \lambda (l), \mu (m),\\
& \nu (n), o (o), \pi (p), \theta (q), \rho (r), \sigma (s), \tau (t), \upsilon (u), \varpi (v), \omega (w), \xi (x), \psi (y), \zeta (z)
}\end{equation}
以下是 script 字母，只有大写有效．所谓大写 $\epsilon$ 其实是花体的 $E$． 
\begin{equation}
\mathcal{A, B, C, D, E, F, G, H, I, J, K, L, M, N, O, P, Q, R, S, T, U, V, W, X, Y, Z}
\end{equation}
另外，电介质常数一律用 $\epsilon$ 而不是 $\varepsilon$．

写单位， 用 \textbackslash Si， 如 $a = 100\Si{m/s^2}$．

% 由于产生各种混乱，以后一律不用文字环绕图片
%\begin{wrapfigure}{r}{0.36\textwidth}
%\vskip-10pt
%\centering
%\includegraphics[width=0.35\textwidth]{./figures/Sample.pdf}
%\caption{A gull}
%\vspace{-15pt}
%\end{wrapfigure}

\subsection{图表}

现在来引用一张图片，图片必须以 eps 以及 pdf 两种格式放在 figures 文件夹中， 并生成 svg 格式放在 littleshi.cn/PhysWiki/online 文件夹中． 生成 svg 时， 字体一律用 outline， 更多选项中保留 7 位小数． 代码中使用 pdf 图片．图片宽度一律用 cm 为单位．
\begin{figure}[ht]
\centering
\includegraphics[width=5cm]{./figures/Sample.pdf}
\caption{例图} \label{Sample_fig1}
\end{figure}
在\autoref{Sample_fig1} 中，label 只能放在 caption 的后面，否则编号会出错．由于图片是浮动的，避免使用“上图”，“下图”等词．

再来看一个表格，如\autoref{Sample_tab1}． 注意标签要放在 caption 后面．
\begin{table}[ht]
\centering
\caption{极限 $\E$ 数值验证}\label{Sample_tab1}
\begin{tabular}{|c|c|c|c|c|c|c|}
\hline
$x$ & ${10^{ - 1}}$ & ${10^{ - 2}}$ & ${10^{ - 3}}$ & ${10^{ - 4}}$ & ${10^{ - 5}}$ & ${10^{ - 6}}$ \\
\hline
$(1 + x)^{1/x}$ & 2.59374 & 2.70481 & 2.71692 & 2.71815 & 2.71827 & 2.71828 \\
\hline
\end{tabular}
\end{table}

下面我们举一个例子并引用

\begin{exam}{名称}\label{Sample_ex1}
在例子中，我们的字体可以自定义，包括公式的字号会保持与内容一致．
\begin{equation}
(a+b)^n = \sum_{i=0}^n C_n^i a^i b^{n-i} \quad (n\text{为整数})
\end{equation}
\end{exam}
引用例子同样使用 \textbackslash autoref， 如\autoref{Sample_ex1}． 

下面是一个习题（\autoref{Sample_exe1}）， 除了标题外其他和例子相同
\begin{exer}{名称}\label{Sample_exe1}
已知 $b$ 和 $c$， 求 $a^2 + b^2 = c^2$．
\end{exer}

\subsection{Matlab 代码}
在行内显示少量代码用 \textbackslash x 命令， 如 \x{plot(x,y)}．
\subsubsection{显示  Command Window 中的代码}
注意方括号如果要出现在行首， 必须用大括号括起来， 否则会出错．
\begin{Command}
>> 1.2/3.4 + (5.6+7.8)*9 -1 \\
ans = 119.9529 \\
>> 1/exp(1) \\
ans = 0.3679 \\
>> exp(-1i*pi)+1 \\
ans = 0
\end{Command}

\subsubsection{显示 m 文件中的代码}

代码必须以  .tex 文件格式放在 code 文件夹中的中， 并用 \textbackslash Code（含文件名）或  \textbackslash code （不含文件名）命令导入正文． 较长的 .tex 代码文件必须含文件名， 文件名需要反映代码的内容而不是与词条 label 同名． 对每个 \textbackslash Code 的 .tex 文件， 需要生成同名的 .m 文件（使用 ANSI 编码）． 代码修改测试以 .tex 文件为准， 把变化后的内容粘贴到 .m 文件中． 这些 .m 文件用于上传到网站上提供下载， 以及用于转换成网页． \textbackslash code 导入的文件命名规则与图片命名相同．

\Code{sample}

%\subsection{文献引用}
%每个词条若有参考文献需要生成与词条文件同名的 bib 文件并放在 bibliography 文件夹中． 文献引用的格式如\cite{PhysWiki}， 在引用了文献的词条最后加上以下两条命令即可生成“ 参考文献” 章节． 如果参考文献有改动， 需要先编译一次 XeLaTeX， 编译一次 BibTeX， 再编译两次 XeLaTeX 即可． 如果文献没有改动则不用编译 BibTeX．

%\bibliography{./bibliography/Sample}
%\bibliographystyle{plain}

\eentry{本书格式规范\upref{Sample}}
\rentry{本书格式规范\upref{Sample}}
